\documentclass[11pt]{article}
\usepackage{amsmath,amssymb}
\usepackage[margin=1in]{geometry}

\newcommand{\HBN}{H_{\mathrm{BN}}}
\newcommand{\MBN}{M_{\mathrm{BN}}}
\newcommand{\XiBN}{\Xi_{\mathrm{BN}}}
\newcommand{\one}{\mathbf 1}
\newcommand{\rhoa}{\rho_a}
\newcommand{\CRE}{\mathrm{CRE}}
\newcommand{\CRCFT}{\mathrm{CRCFT}}

\title{Step 193: Beurling--Nyman Carrier Pivot}
\date{}

\begin{document}
\maketitle

\section*{Verdict}

\[
\boxed{\texttt{V\_BN\_CRCFT\_TE}}.
\]

The Beurling--Nyman carrier is a new classically-RH-equivalent carrier with a different attack shape from Burnol/Sonine.  It nevertheless lands in \(\CRCFT\)-TE: the closure condition is exactly target-equivalent to RH.

\section*{Inherited Audit}

The inherited records contain only sparse Beurling--Nyman references: a Step 93 warning that small carriers are underpowered against the de Branges/Baez--Duarte/Bombieri wall, and the Step 193 manager rationale in the cascade map.  No accepted BN carrier package was inherited.  This step therefore declares the carrier directly from the classical Beurling--Nyman literature.

\section*{Carrier Declaration}

Let
\[
\HBN=L^2(0,1)
\]
with Lebesgue measure.  For \(0<a\le 1\), define
\[
\rho_a(t)=\left\{\frac{a}{t}\right\}-a\left\{\frac{1}{t}\right\},
\qquad 0<t<1,
\]
where \(\{x\}\) is the fractional part.  Let
\[
\MBN=\overline{\operatorname{span}\{\rho_a:0<a\le1\}}^{\,L^2(0,1)}.
\]
The parent residual is
\[
\XiBN=(I-P_{\MBN})\one,
\]
and its scalar defect is
\[
\delta_{\mathrm{BN}}^2
=
\|\XiBN\|_2^2
=
\inf_{f\in\MBN}\|\one-f\|_2^2.
\]

\section*{Classical Criterion}

The Beurling--Nyman theorem states:
\[
\XiBN=0
\quad\Longleftrightarrow\quad
\one\in\MBN
\quad\Longleftrightarrow\quad
\mathrm{RH}.
\]
Thus the carrier is \(\CRE\).  Baez--Duarte's strengthening permits restricting the approximation family to countable arithmetic parameter families, in particular integer-dilation versions of the criterion.

\section*{Gram Matrix Subinterface}

For a finite set \(A=\{a_1,\dots,a_N\}\subset(0,1]\), set
\[
G_A(i,j)=\langle \rho_{a_i},\rho_{a_j}\rangle_{L^2(0,1)}
=
\int_0^1
\left(\left\{\frac{a_i}{t}\right\}-a_i\left\{\frac1t\right\}\right)
\left(\left\{\frac{a_j}{t}\right\}-a_j\left\{\frac1t\right\}\right)
\,dt.
\]
Also
\[
b_A(i)=\langle \one,\rho_{a_i}\rangle
=a_i\log(1/a_i).
\]
When \(G_A\) is invertible, the finite-window distance is
\[
\delta_A^2
=
1-b_A^*G_A^{-1}b_A,
\]
and in general \(G_A^{-1}\) is replaced by the Moore--Penrose inverse \(G_A^\dagger\).

The Mellin transform identity behind the criterion is
\[
\int_0^1 \rho_a(t)t^{s-1}\,dt
=
\frac{a-a^s}{s}\zeta(s),
\]
with the usual meromorphic interpretation at \(s=1\).  This identity explains why the closure of \(\one\) by the \(\rho_a\)'s is a zeta-zero statement.

The matrix elements are computable: one may partition \((0,1)\) at the breakpoints \(a_i/n\), \(a_j/m\), \(1/n\), or use the arithmetic formulas developed in the Baez--Duarte literature.  Hence this is not a missing-kernel CTMT blocker.

\section*{CRCFT Classification}

The natural BN finite-window computations are matrix computations, but the matrix entries themselves are available.  The hard assertion is
\[
\inf_A \delta_A^2=0
\]
over an exhausting family of parameters.  By the Beurling--Nyman theorem this is equivalent to RH.  Therefore the framework mode is:
\[
\boxed{\CRCFT\text{-TE}.}
\]
The Gram system is a computational subinterface, not a CTMT-stuck terminal missing record.

\section*{Dichotomy Update}

The BN carrier is the sixth CRE carrier tested.  It confirms the Dichotomy pattern: a fresh CRE carrier with a different shape still falls into a CRCFT foreclosure mode, here target-equivalence.

\section*{Nonclaim}

This step does not establish RH and does not close \(\XiBN\).  It only classifies the carrier and records the finite Gram-matrix structure.

\end{document}

